\section{Preliminaries and conventions}
\label{sec:prelim}

\begin{convention}\label{conv:standing}
Throughout, an \emph{algebra} means a unital real Banach algebra $A$ whose norm
is submultiplicative with $\norm{1}=1$.  Homomorphisms are unital and
$\RR$-linear.  We write $\Zc(A)$ for the centre, $\rad(A)$ for the Jacobson
radical, $[a,z]=az-za$ and $\ad_a=[a,-]$.  We allow the \emph{degenerate
algebra} $0$, in which $1=0$; see Convention~\ref{conv:degenerate}.
\end{convention}

\subsection{Complexification and spectrum}

For a real algebra $A$ let $A_\CC=A\otimes_\RR\CC$ carry any of the usual
equivalent norms, e.g.\
$\norm{x+iy}=\sup_{\theta}\norm{x\cos\theta-y\sin\theta}$, which is
submultiplicative and restricts to the given norm on $A$.  The \emph{spectrum}
of $a\in A$ is
\[
  \spec(a)=\{\lambda\in\CC:\lambda1-a\ \text{is not invertible in}\ A_\CC\},
\]
a nonempty compact subset of $\CC$ stable under complex conjugation, and the
spectral radius is $r(a)=\max\{\abs{\lambda}:\lambda\in\spec(a)\}$.  Beurling's
formula $r(a)=\lim_n\norm{a^n}^{1/n}$ holds verbatim, since the norms on $A$ and
$A_\CC$ are equivalent within a factor of $2$ and the limit is insensitive to
that.  A \emph{character} of $A$ is a nonzero homomorphism $\chi:A\to\CC$ of
real algebras with $\chi(1)=1$; every character is automatically contractive,
and $\abs{\chi(a)}\le r(a)$.

\begin{remark}
Nothing forces a noncommutative $A$ to have any characters at all: $M_2(\RR)$
has none, since a character would kill all commutators and
$M_2(\RR)=\RR1\oplus[M_2(\RR),M_2(\RR)]$, while the resulting functional
$a\mapsto\tfrac12\tr(a)$ is not multiplicative.  Statements below that quantify
over characters are therefore vacuous on such algebras; this is a feature, not a
gap.  The geometric content is carried by the centre instead (\S\ref{sec:base}).
\end{remark}

\subsection{Real Gelfand theory}

We use the real-scalar Gelfand theory in the form developed by Kulkarni and
Limaye \cite{KulkarniLimaye}.  For commutative $A$ the character space
$\Omega(A)$ of $A_\CC$ is compact Hausdorff and carries the involution
$\chi\mapsto\bar\chi$, $\bar\chi(a)=\overline{\chi(a)}$; the Gelfand transform
\[
  \Gamma:A\longrightarrow C(\Omega(A),\CC),\qquad \Gamma(a)(\chi)=\chi(a)
\]
has image inside the \emph{real function algebra} of conjugation-equivariant
functions, $\norm{\Gamma(a)}_\infty=r(a)$, and $\ker\Gamma=\rad(A)$.  The
involution is trivial --- so that the image consists of real-valued functions
--- exactly when every character is real-valued.  We write $\CompHaus$ for the
category of compact Hausdorff spaces and continuous maps.

\begin{theorem}[Real Stone--Weierstrass]\label{thm:stone-weierstrass}
A closed unital subalgebra of $C(X,\RR)$, $X$ compact Hausdorff, which separates
the points of $X$ is all of $C(X,\RR)$.
\end{theorem}

\subsection{Two classical inputs}

We shall use the following two results as black boxes; both are standard, and
both are the reason the rigid core of \S\ref{sec:rigid} is as rigid as it is.

\begin{theorem}[Hirschfeld--\.Zelazko \cite{HirschfeldZelazko}]
\label{thm:hz}
Let $B$ be a complex Banach algebra whose spectral radius is submultiplicative,
$r(xy)\le r(x)r(y)$ for all $x,y$.  Then $B/\rad(B)$ is commutative.
\end{theorem}

\begin{theorem}[Stampfli \cite{Stampfli}]\label{thm:stampfli}
For $a\in B(H)$, $H$ a complex Hilbert space,
$\norm{\ad_a}=2\dist(a,\CC1)$.
\end{theorem}

Theorem~\ref{thm:stampfli} is stated over $\CC$; the real analogue we use ---
$\norm{\ad_a}=2\dist(a,\RR1)$ for $a\in M_n(\RR)$ acting on $\RR^n$ --- is
recorded as Proposition~\ref{prop:real-stampfli} and is, in the present
treatment, verified numerically rather than proved.
